%%=============================================================================
%% Conclusie en reflectie
%%=============================================================================

\chapter{\IfLanguageName{dutch}{Conclusie en reflectie}{Conclusion and reflection}}%
\label{ch:reflectie}
TODO!
%% Richtlengte: 1 tot 2 bladzijden.
Het graduaatsproject vormde een concrete uitbreiding van het Vesalius-platform op basis van twee bestaande JIRA-tickets. Het project combineerde analyse van een bestaande codebase met het ontwerpen, implementeren en testen van nieuwe functionaliteit binnen een professionele ontwikkelomgeving.


\section{\IfLanguageName{dutch}{Conclusie}{Conclusion}}%
\label{sec:conclusie}
TODO!
%% Geef een duidelijk antwoord op de doelstelling uit je inleiding. Is ze
%% gehaald, gedeeltelijk gehaald, of niet? Wat levert je werk op voor het
%% bedrijf?
%%
%% Verwijs terug naar wat je in Hoofdstuk 1 beloofde. Een conclusie die geen
%% enkele band heeft met de doelstelling, leest als een ander project.
De doelstelling van het project was om automatisch gegenereerde output binnen Vesalius.ai gerichter te kunnen afstemmen op de gebruiker en doelgroep.
Voor documenttemplates werd hiervoor een mechanisme uitgewerkt waarmee templates optioneel aan specifieke artsen kunnen worden gekoppeld. Hierdoor kan automatische generatie gericht worden uitgevoerd zonder de bestaande zichtbaarheid van templates of de mogelijkheid tot manuele generatie te wijzigen.
Voor Scribe-consultatemplates werd een doelgroepinstelling voorzien waarmee onderscheid kan worden gemaakt tussen patiëntgerichte en artsgerichte output.
De belangrijkste meerwaarde voor Vesalius.ai is dat automatische output relevanter kan worden gemaakt voor de gebruiker. Tegelijk blijft de bestaande werking behouden wanneer de nieuwe configuratiemogelijkheden niet worden gebruikt. Hierdoor vormt de uitbreiding geen volledige herwerking van het bestaande systeem, maar een gerichte verbetering binnen de bestaande architectuur.

\section{\IfLanguageName{dutch}{Wat ik geleerd heb}{What I learned}}%
\label{sec:geleerd}
TODO!
%% Zowel technisch als over jezelf als werkende professional. Wees concreet:
%% "ik heb leren werken met X" zegt weinig, "ik onderschatte hoeveel tijd het
%% opzetten van X kost, en ik plan dat voortaan apart in" zegt iets.
Een belangrijk leerpunt was het werken binnen een bestaande professionele codebase. In een eigen schoolproject kan een technische oplossing vanaf het begin worden ontworpen. Binnen een bestaand product moet daarentegen eerst worden onderzocht waarom de huidige oplossing op een bepaalde manier is opgebouwd en welke bestaande functionaliteit door een wijziging kan worden beïnvloed.
Daarnaast heb ik geleerd dat de functionele beschrijving van een JIRA-ticket niet automatisch gelijkstaat aan een technische oplossing. De acceptatiecriteria beschrijven wat de gebruiker moet kunnen, maar het is aan de ontwikkelaar om te bepalen hoe dit binnen de bestaande architectuur gerealiseerd wordt.
Ook het belang van backwards compatibility werd duidelijk. Een nieuwe instelling toevoegen is niet voldoende. Er moet eveneens worden gegarandeerd dat bestaande configuraties zonder de nieuwe instelling zich blijven gedragen zoals voordien.
Op professioneel vlak heb ik verder ervaring opgedaan met Git, Jira, code review, een bestaande CI/CD-omgeving, Docker en de samenwerking binnen een developmentteam. Hierdoor kreeg ik een beter inzicht in het verschil tussen een schoolproject en software die deel uitmaakt van een effectief gebruikt product.

\section{\IfLanguageName{dutch}{Wat ik anders zou doen}{What I would do differently}}%
\label{sec:anders}
TODO!
%% De vraag die vrijwel elke jury stelt. Wie ze vooraf beantwoord heeft, komt
%% sterker over dan wie ter plaatse begint na te denken.
Een volgende keer zou ik nog vroeger een volledige analyse van de bestaande functionaliteit uitvoeren voordat ik met de implementatie begin. Vooral bij een bestaand systeem kan een groot deel van het werk bestaan uit het begrijpen van bestaande relaties, businesslogica en uitzonderingen.
Daarnaast zou ik de acceptatiecriteria vanaf het begin vertalen naar concrete testscenario's. Hierdoor kan tijdens de ontwikkeling sneller worden gecontroleerd of ieder onderdeel van de oorspronkelijke vraag effectief afgedekt is.
Ook zou ik ontwerpbeslissingen en openstaande vragen systematischer documenteren tijdens de ontwikkeling. Dit maakt het eenvoudiger om later te verklaren waarom een bepaalde oplossing gekozen werd en voorkomt dat belangrijke beslissingen enkel in het geheugen of in gesprekken met teamleden blijven bestaan.

\section{\IfLanguageName{dutch}{Verder werk}{Future work}}%
\label{sec:verder-werk}
TODO!
%% Wat zou een volgende stap zijn voor het bedrijf? Wat blijft er liggen, en
%% wat is daarvoor nodig?
Een mogelijke volgende stap is het verder uitbreiden van de templatefunctionaliteit naar andere templatecategorieën wanneer daar vanuit het product nood aan bestaat. Daarbij moet telkens eerst worden nagegaan of de betrokken templatecategorie technisch en functioneel op dezelfde manier werkt.
Een afzonderlijk toekomstig project is het WYSIWYG-beheer van e-mails. Daarbij zou niet alleen een gebruiksvriendelijke editor moeten worden voorzien, maar ook een betrouwbare manier om compatibele e-mail-HTML te genereren. Omdat verschillende e-mailclients HTML en CSS op verschillende manieren interpreteren, vereist dit een afzonderlijke technische analyse en teststrategie.
De functionaliteit uit dit graduaatsproject kan daarnaast verder worden uitgebreid wanneer het team nieuwe vereisten rond doelgroepen, gebruikersspecifieke configuratie of automatische output toevoegt.